\clearpage
\section*{September 8, 2026: Carry individual locations through to boundary distances}
\addcontentsline{toc}{section}{September 8, 2026: Carry individual locations through to boundary distances}
\textit{Agent-drafted entry.}

The approved same-property rule now locates 2,052 of the 2,078 previously flagged residential records. The comparison uses an exact parcel point from an assessment year reporting the selected units, floor area, land area and construction year, with the existing committed year corrections respected. It chooses the latest qualifying saved point, including older assessments when the property is absent or changed in 2025. The point must fall within the historical site; a tolerance of 0.001 foot handles coordinate rounding. This locates the property without replacing its reported land area with a polygon area.

Nineteen further records already had exact Chicago address points accepted upstream. The geography script now carries these points through instead of replacing them with the historical parcel center. This includes the previously approved South Troy address correction. Altogether, 2,071 of the original 2,078 locations are resolved under these rules. Of the 69 initial exceptions, 62 are resolved and seven remain outside the rules. The rules belong in the existing geography producer and use no new observation-specific production decisions.

Relative to the reconstructed parcel centers, eleven of the 2,071 records enter the 500-foot band and 22 leave. All 788 records matching the frozen paper input by exact project identifier keep their previous membership in that band. The originally tested 2,009 points reproduce the audit's proposed distances exactly. These checks do not establish unchanged regression coefficients. The paper input remains unchanged and no density regressions were run.

Seven remaining location cases have identifiable data limitations. Greenview and Lowe have later floor-area changes; Indiana has a later construction-year change. Avers, 34th Street unit D, Sangamon and Oakenwald no longer have residential assessments under their recorded parcel numbers. Avers has both a 2006 new-building permit and a 2014 demolition permit, so later vacancy does not negate earlier construction. Exempt ownership is a plausible explanation at 34th Street, where the Chicago Housing Authority also lists the address. The other classification changes do not establish demolition or a replacement parcel identity.

Location and construction eligibility must remain separate. Wabash's location is now established, but its 2009 permit explicitly remodels an existing building, and its later assessment reports an 1890 year instead of 2010. Exclusion of that purported construction is recommended. Indiana's shift from 2021 to 1886 likewise warrants holding it out unless a replacement building is documented. Greenview's selected 2012 year precedes its June 2013 new-building permit and requires completion evidence. These three concerns are preserved in the consolidated questions; new exclusions or guessed completion years were not applied in this change. Later floor-area revisions are recorded without automatically replacing the selected measurements.

\paragraph{Reproduction.} The existing preferred-geography producer and boundary-scope task generate the revised points and distances. The release audit's \texttt{remaining\_individual\_location\_findings.csv} follows the fixed 69-record cohort; \texttt{remaining\_individual\_locations.md} records exact permit identifiers and supporting sources. Make builds, standard reports, distance comparisons, and unchanged-build checks verify this bounded change. Final construction assembly remains unfinished.
